\part{Memory management}

\chapter{The physical memory manager}

\begin{goals}
  \item Why memory is handed out in fixed-size chunks
  \item What a bitmap allocator is and how it works
  \item How the kernel protects itself from its own allocator
  \item The rounding bug that makes it fail, and why it is so hard to spot
\end{goals}

\section{Pages and frames}

The first decision any memory manager makes is the size of the unit it deals in. lytlnyblOS, like
every x86 system, uses \hi{4096 bytes}. The hardware forces this: the paging unit we will meet
in the next chapter works in 4 KB granules and nothing else.

The vocabulary matters and people mix it up constantly:

\begin{center}
\begin{tabularx}{\textwidth}{@{}l X@{}}
\toprule
\textbf{Term} & \textbf{Meaning} \\
\midrule
\term{Frame} & A 4 KB chunk of \emph{physical} RAM. A real, existing piece of hardware \\
\term{Page}  & A 4 KB chunk of \emph{virtual} address space. A name that may or may not refer to anything \\
\bottomrule
\end{tabularx}
\end{center}

The physical manager in this chapter deals only with \emph{frames}. It answers one question, and
nothing else: \emph{which 4 KB pieces of real RAM are in use?}

\begin{analogy}
Think of a hotel. The frames are the rooms: they physically exist, there is a fixed number of them,
and each is either occupied or free. The physical manager is the receptionist with the occupancy
board. It does not know who is in which room, what they are doing, or why --- only occupied or
free.

The next chapter introduces room numbers that guests see, which need not match the real ones. That
is virtual memory.
\end{analogy}

\section{A bitmap}

With 128 MB of RAM there are $128 \times 1024 / 4 = 32{,}768$ frames. How do you track 32,768
yes/no facts? With 32,768 bits, that is 4 KB total. That is a \term{bitmap}: one bit per frame, 1 =
used, 0 = free.

\filetag{lytlnybl/memory/pmm.c}
\begin{lstlisting}[style=c]
uint32_t* bitmap = (uint32_t*)BITMAP_BASE;   /* BITMAP_BASE is 0x100000 */

void bitmap_set_frame(uint32_t frame_index) {
    bitmap[frame_index / 32] |= (1 << (frame_index % 32));
}

void bitmap_clear_frame(uint32_t frame_index) {
    bitmap[frame_index / 32] &= ~(1 << (frame_index % 32));
}
\end{lstlisting}

The bitmap is an array of \cod{uint32\_t}, so each element holds 32 frames. To reach frame $i$:

\begin{itemize}
  \item \cod{i / 32} says which array element,
  \item \cod{i \% 32} says which bit within it,
  \item \cod{1 << (i \% 32)} builds a mask with only that bit set.
\end{itemize}

\begin{center}
\begin{tikzpicture}[font=\ttfamily\scriptsize]
  \foreach \i in {0,...,7} {
    \draw[draw=grayborder,fill=graylight] (\i*0.72,0) rectangle (\i*0.72+0.66,0.66);
  }
  \node at (0.33,0.33) {1}; \node at (1.05,0.33) {1}; \node at (1.77,0.33) {0};
  \node at (2.49,0.33) {0}; \node at (3.21,0.33) {1}; \node at (3.93,0.33) {0};
  \node at (4.65,0.33) {0}; \node at (5.37,0.33) {0};
  \node[anchor=west,font=\sffamily\scriptsize,text=gray] at (6.0,0.33) {\ldots\ 32 bits = one uint32\_t};
  \foreach \i/\l in {0/f0,1/f1,2/f2,3/f3,4/f4,5/f5,6/f6,7/f7}
    \node[font=\ttfamily\tiny,text=gray] at (\i*0.72+0.33,-0.3) {\l};
  \draw[decorate,decoration={brace,amplitude=3pt},green] (1.44,-0.55) -- (3.15,-0.55)
    node[midway,below,font=\sffamily\tiny,text=green] {free};
\end{tikzpicture}
\end{center}

\section{Allocating and freeing}

\filetag{lytlnybl/memory/pmm.c}
\begin{lstlisting}[style=c]
void* alloc_frame() {
    uint32_t frame_i;
    for (frame_i = 0; frame_i < total_frames; frame_i++) {
        uint32_t is_reserved = (bitmap[frame_i / 32] & (1 << (frame_i % 32)));
        if (!is_reserved) {
            bitmap_set_frame(frame_i);
            return (void *)(frame_i * FRAME_SIZE);
        }
    }
    return NULL;    /* out of memory */
}

void free_frame(void* frame_address) {
    uint32_t frame_i = (uint32_t)frame_address / FRAME_SIZE;
    bitmap_clear_frame(frame_i);
}
\end{lstlisting}

This is \term{first fit}: scan from the beginning and take the first free frame. Simple, correct,
and $O(n)$ --- with 32,768 frames the loop can run 32,768 times for a single allocation.

\begin{warn}{Why it is fast enough here, and would not be in general}
Because allocations are rare (a few at boot, a handful per process) and because free frames tend to
be near the start, the loop usually exits after a few iterations. A real system keeps a free list, a
pointer to the last-found position, or a multi-level bitmap. Worth knowing the limit exists, not
worth fixing here.
\end{warn}

\begin{danger}{\cod{free\_frame} validates nothing}
It does not check that the address is page-aligned, that it is within range, or that the frame was
actually allocated. \cod{free\_frame((void*)0x1001)} silently clears the bit for frame 1, and
\cod{free\_frame} called twice on the same frame hands it out to two owners. In a kernel, that is
the seed of arbitrarily strange bugs appearing arbitrarily far away.
\end{danger}

\section{Initialisation: deciding what is off limits}

This is the delicate part. The manager starts by marking \emph{everything} reserved and then frees
only what the BIOS map says is usable:

\filetag{lytlnybl/memory/pmm.c}
\begin{lstlisting}[style=c]
/* mark every region reserved, for safety */
for (uint32_t i = 0; i < bitmap_entries; i++) {
    bitmap[i] = 0xFFFFFFFF;
}

/* then free the usable regions */
for (uint32_t i = 0; i < map_entry_count; i++) {
    if (memory_map[i].type == 1) {                /* usable RAM */
        uint64_t starting_frame = memory_map[i].base / FRAME_SIZE;
        uint64_t frame_length   = memory_map[i].length / FRAME_SIZE;
        for (uint32_t j = starting_frame; j < starting_frame + frame_length; j++) {
            bitmap_clear_frame(j);
        }
    }
}
\end{lstlisting}

\begin{concept}{Deny by default}
Starting with everything reserved and opening up only what is explicitly known to be safe is a
security principle worth internalising. The opposite --- start everything free and mark the bad
parts --- fails catastrophically the moment you forget one region, and it fails silently.
\end{concept}

Then come the specific protections: the BIOS area, the kernel itself, the bitmap, the boot sector
and video memory.

\begin{lstlisting}[style=c]
bitmap_set_frame(0);                              /* BIOS data */

uint32_t kernel_frame_index = (uint32_t)&_kernel_start / FRAME_SIZE;
uint32_t kernel_frame_end   = (uint32_t)&_kernel_end   / FRAME_SIZE;
for (; kernel_frame_index < kernel_frame_end; kernel_frame_index++) {
    bitmap_set_frame(kernel_frame_index);         /* the kernel */
}

uint32_t bitmap_frame_start = BITMAP_BASE / FRAME_SIZE;
uint32_t bitmap_size_bytes  = bitmap_entries * sizeof(uint32_t);
uint32_t bitmap_frame_count = (bitmap_size_bytes + FRAME_SIZE - 1) / FRAME_SIZE;
for (uint32_t i = 0; i < bitmap_frame_count; i++) {
    bitmap_set_frame(bitmap_frame_start + i);     /* the bitmap itself */
}

bitmap_set_frame(0x7C00 / FRAME_SIZE);            /* the boot sector */

uint32_t video_frame_index = 0xA0000 / FRAME_SIZE;
uint32_t video_frame_end   = 0xFFFFF / FRAME_SIZE;
for (; video_frame_index < video_frame_end; video_frame_index++) {
    bitmap_set_frame(video_frame_index);          /* video memory */
}
\end{lstlisting}

\section{The rounding bug}

Look carefully at two lines above. The bitmap-size calculation uses the round-up idiom:

\begin{lstlisting}[style=c]
uint32_t bitmap_frame_count = (bitmap_size_bytes + FRAME_SIZE - 1) / FRAME_SIZE;
\end{lstlisting}

The kernel-end calculation does not:

\begin{lstlisting}[style=c]
uint32_t kernel_frame_end = (uint32_t)&_kernel_end / FRAME_SIZE;
\end{lstlisting}

\begin{danger}{The consequence, with real numbers}
In one build of this project, \cod{\_kernel\_end} is \hexa{13468}.

$$\hexa{13468} \div \hexa{1000} = 19.27 \rightarrow 19 \text{ (integer division truncates)}$$

The loop condition is \cod{index < 19}, so it protects frames 9 through \textbf{18}: addresses
\hexa{9000} to \hexa{12FFF}. \hi{Frame 19 --- \hexa{13000} to \hexa{13FFF} --- is left free},
even though the last 1,128 bytes of the kernel live in it.

What lives in those 1,128 bytes? In that build: \cod{idtr}, \cod{terminal}, \cod{keymap},
\cod{kernel\_directory}, \cod{current\_directory}, \cod{heap\_start\_head}, \cod{open\_files},
\cod{process\_head}, \cod{current\_process} and the \cod{tss}. Essentially every critical global in
the kernel.

So \cod{alloc\_frame()} eventually returns \hexa{13000}, the loader copies a user program's code
into it, and the kernel's own state is destroyed. \cod{terminal.width} and \cod{terminal.height}
become whatever bytes of the shell binary landed there --- possibly zero. And a zero \cod{width}
means the divide-by-zero cascade from Chapter 20.

The fix is one line, and it is \hi{now applied in the source tree}:
\begin{lstlisting}[style=c]
uint32_t kernel_frame_end = ((uint32_t)&_kernel_end + FRAME_SIZE - 1) / FRAME_SIZE;
\end{lstlisting}
\end{danger}

\begin{warn}{Why this bug is so nasty}
It is \hi{latent}. Whether it bites depends on exactly where \cod{\_kernel\_end} falls and
which globals end up in the last partial page --- which changes every time you add or remove code.
A build that works perfectly can break because you added a comment that changed a string constant's
length.

That is the worst class of bug: one where ``it works on my machine'' and ``it works on the same
machine yesterday'' are both true, and both meaningless. Part XIII documents how it was actually
found.
\end{warn}

\begin{recipe}{asking for and returning physical frames}{ch25-physical-memory}
Allocates three frames, prints their physical addresses, fills one with \cod{memset} to prove it is
real memory, frees the middle one and allocates again --- getting the same frame straight back,
because the allocator is first fit.

The addresses it prints are worth reading: the third frame is not adjacent to the second, because
the frame holding the boot sector at \hexa{7C00} is protected and gets skipped.
\end{recipe}


\exercise{The region between \hexa{1000} and \hexa{8FFF} is marked usable by the BIOS map and is
never explicitly protected. The E820 memory map itself was stored at \hexa{5000}. Is that a bug?
Justify your answer by considering \emph{when} the map is read for the last time.}

% ============================================================
\chapter{Virtual memory, slowly}

\begin{goals}
  \item Understand, with no prior knowledge, what virtual memory is and what problem it solves
  \item Follow an address translation by hand, with real numbers
  \item Understand why there are two levels of table and not one
  \item Know what every bit of a page table entry means
\end{goals}

This is the hardest chapter in the course and the one everything else rests on. It is written to be
read slowly. If a section does not click, do not push on --- reread it. Nothing later makes sense
without this.

\section{The problem, told as a story}

Suppose we have two programs and we want to run both.

Program A was compiled to live at address \hexa{400000}. Its code says things like ``jump to
\hexa{4001A0}'' and ``the variable \cod{counter} is at \hexa{400C20}''. Those addresses are baked
into the machine code at compile time.

Program B was compiled the same way --- because the compiler always puts programs at
\hexa{400000}. Its code also says ``jump to \hexa{4001A0}''.

\begin{center}
\begin{tikzpicture}[font=\sffamily\scriptsize]
  \node[draw=green,fill=greenlight,rounded corners=3pt,minimum width=3.2cm,minimum height=1.6cm,align=center] (a) at (0,0)
    {\textbf{Program A}\\{\tiny wants to live at}\\ \texttt{0x400000}};
  \node[draw=purple,fill=purplelight,rounded corners=3pt,minimum width=3.2cm,minimum height=1.6cm,align=center] (b) at (0,-2.2)
    {\textbf{Program B}\\{\tiny wants to live at}\\ \texttt{0x400000}};
  \node[draw=red,fill=redlight,rounded corners=3pt,minimum width=3.4cm,minimum height=3.8cm,align=center] (r) at (6,-1.1)
    {\textbf{Physical RAM}\\[4pt]{\tiny there is only ONE}\\{\tiny address 0x400000}};
  \draw[-{Stealth[length=2.5mm]},green,thick] (a) -- (r);
  \draw[-{Stealth[length=2.5mm]},purple,thick] (b) -- (r);
  \node[text=red,font=\sffamily\bfseries\small] at (3.1,-1.1) {conflict};
\end{tikzpicture}
\end{center}

There is only one physical address \hexa{400000}. Both cannot have it.

\subsection{Three bad solutions}

\begin{description}
  \item[Recompile every program for a free address] Then you cannot run the same binary twice, and
    you must know at compile time what else will be running. Absurd.
  \item[Relocate the code at load time] Rewrite every address in the binary as you load it. This
    was actually done historically. It is slow, error-prone, and requires the binary to carry a
    table of every address to patch.
  \item[Ask programs to be polite and avoid each other] This requires trusting every program on the
    machine. One buggy program takes down everything. This is real mode, and it is why real mode is
    useless.
\end{description}

\subsection{The good solution}

Give each program \hi{its own private set of addresses}, and have the hardware translate
those private addresses into real ones on every single access.

Program A says \hexa{400000} and the hardware quietly turns it into physical \hexa{800000}.
Program B says \hexa{400000} and the hardware quietly turns it into physical \hexa{C00000}.
Neither program knows. Neither program can find out. Neither program can reach the other's memory,
because \emph{there is no address it can write down} that leads there.

\begin{concept}{The central idea of this chapter}
A \term{virtual address} is a name. A \term{physical address} is a place. Virtual memory is a
\hi{dictionary from names to places}, and each program gets its own dictionary.

Two programs can use the same name because their dictionaries are different. A program cannot reach
another's memory because its dictionary has no entry leading there. Isolation is not enforced by
checking --- it is enforced by \emph{the absence of a translation}.
\end{concept}

\begin{analogy}[Back to the hotel]
Every guest is handed a map on which their room is labelled ``Room 1''. Guest A's Room 1 is really
room 314; guest B's Room 1 is really room 502. Both can say ``I'm going to Room 1'' and both arrive
somewhere different, without confusion and without ever meeting.

If guest A tries to walk to ``Room 7'' and their map has no Room 7, they do not end up in someone
else's room --- they end up in front of a wall. The hotel manager (the kernel) gets a call about it.
That call is a \emph{page fault}.
\end{analogy}

\section{Why a dictionary of every address is impossible}

The obvious implementation would be a table with one entry per address: 4 billion entries. At 4
bytes each that is 16 GB of table, to manage 4 GB of memory. Obviously absurd.

The fix is to translate in \hi{blocks}, not individual bytes. If we translate in 4 KB blocks,
we need only $4\,\text{GB} / 4\,\text{KB} = 1{,}048{,}576$ entries. At 4 bytes each: 4 MB of table.

Better --- but still 4 MB \emph{per process}. Ten processes would be 40 MB of pure bookkeeping,
most of it describing addresses nobody uses.

\subsection{The trick: split the table in two}

Instead of one table of 1,048,576 entries, use:

\begin{itemize}
  \item one small table of \hi{1024} entries, called the \term{page directory};
  \item where each entry points to another table of \hi{1024} entries, called a \term{page
        table};
  \item where each of \emph{those} entries describes one 4 KB page.
\end{itemize}

$1024 \times 1024 = 1{,}048{,}576$. The coverage is identical. But the saving is enormous, because
\hi{a directory entry can be empty}. A process that uses 8 KB of code and 4 KB of stack needs
a directory (4 KB) and one or two page tables (4--8 KB) --- about 12 KB, not 4 MB.

\begin{center}
\begin{tikzpicture}[font=\sffamily\scriptsize]
  \node[draw=accent,fill=accentlight,rounded corners=3pt,minimum width=2.7cm,minimum height=3.4cm,align=center] (d) at (0,0)
    {\textbf{Page directory}\\[2pt]{\tiny 1024 entries}\\{\tiny 4 KB}};
  \node[draw=green,fill=greenlight,rounded corners=3pt,minimum width=2.7cm,minimum height=1.5cm,align=center] (t1) at (5,1.2)
    {\textbf{Page table 0}\\{\tiny 1024 entries}};
  \node[draw=green,fill=greenlight,rounded corners=3pt,minimum width=2.7cm,minimum height=1.5cm,align=center] (t2) at (5,-1.2)
    {\textbf{Page table 767}\\{\tiny 1024 entries}};
  \node[text=gray,font=\sffamily\tiny,align=center] at (5,0) {(the other 1022\\entries are empty)};

  \node[draw=amber,fill=amberlight,rounded corners=3pt,minimum width=2.4cm,minimum height=0.8cm,align=center] (p1) at (9.8,1.7) {4 KB page};
  \node[draw=amber,fill=amberlight,rounded corners=3pt,minimum width=2.4cm,minimum height=0.8cm,align=center] (p2) at (9.8,0.7) {4 KB page};
  \node[draw=amber,fill=amberlight,rounded corners=3pt,minimum width=2.4cm,minimum height=0.8cm,align=center] (p3) at (9.8,-1.2) {4 KB page};

  \draw[-{Stealth[length=2mm]},accent] (d) -- (t1);
  \draw[-{Stealth[length=2mm]},accent] (d) -- (t2);
  \draw[-{Stealth[length=2mm]},green] (t1) -- (p1);
  \draw[-{Stealth[length=2mm]},green] (t1) -- (p2);
  \draw[-{Stealth[length=2mm]},green] (t2) -- (p3);
\end{tikzpicture}
\end{center}

\begin{concept}{Where 1024 and 4096 come from}
These numbers are not arbitrary; they fall out of the arithmetic:
\begin{itemize}
  \item A page is 4096 bytes. To address a byte inside it you need
        $\log_2 4096 = \hi{12 bits}$.
  \item That leaves $32 - 12 = 20$ bits to name the page.
  \item Split those 20 evenly: 10 bits and 10 bits. And
        $2^{10} = \hi{1024}$ --- the size of each table.
  \item Each entry is 4 bytes, so a table of 1024 entries is exactly
        $1024 \times 4 = \hi{4096}$ bytes: one page. The tables fit in the very
        thing they describe.
\end{itemize}
The whole design closes on itself. That elegance is why it has survived since 1985.
\end{concept}

\section{Splitting an address}

Every 32-bit virtual address is read by the hardware as three fields:

\begin{center}
\begin{tikzpicture}[font=\sffamily\scriptsize,x=1cm]
  \draw[draw=grayborder,fill=accentlight] (0,0) rectangle (4.0,1.0);
  \node[align=center] at (2.0,0.5) {\textbf{directory index}\\10 bits};
  \draw[draw=grayborder,fill=greenlight] (4.0,0) rectangle (8.0,1.0);
  \node[align=center] at (6.0,0.5) {\textbf{table index}\\10 bits};
  \draw[draw=grayborder,fill=amberlight] (8.0,0) rectangle (12.8,1.0);
  \node[align=center] at (10.4,0.5) {\textbf{offset within the page}\\12 bits};

  \node[font=\ttfamily\tiny,text=gray] at (0.5,-0.3) {bit 31};
  \node[font=\ttfamily\tiny,text=gray] at (3.6,-0.3) {22};
  \node[font=\ttfamily\tiny,text=gray] at (4.4,-0.3) {21};
  \node[font=\ttfamily\tiny,text=gray] at (7.6,-0.3) {12};
  \node[font=\ttfamily\tiny,text=gray] at (8.5,-0.3) {11};
  \node[font=\ttfamily\tiny,text=gray] at (12.3,-0.3) {0};
\end{tikzpicture}
\end{center}

And the code that does the splitting is exactly the idiom from Chapter 2:

\filetag{lytlnybl/memory/vmm.c}
\begin{lstlisting}[style=c]
uint16_t dir_index   = (virtual_address >> 22);
uint16_t table_index = (virtual_address >> 12 & 0x3FF);
/* the offset is virtual_address & 0xFFF */
\end{lstlisting}

\subsection{A worked example, digit by digit}

Let us translate the virtual address \hexa{00401A2C} completely by hand. This is the kind of address
a user program actually uses in this system.

\textbf{Step 1 --- write it in binary.}

\hexa{00401A2C} is:
\begin{center}
\cod{0000 0000 0100 0000 0001 1010 0010 1100}
\end{center}

\textbf{Step 2 --- cut it into 10, 10 and 12 bits.}

\begin{center}
\begin{tikzpicture}[font=\ttfamily\small,x=1cm]
  \node[anchor=west,text=accent] at (0,0.6) {0000000001};
  \node[anchor=west,text=green]  at (3.4,0.6) {0000000110};
  \node[anchor=west,text=amber]  at (6.8,0.6) {100000101100};
  \node[anchor=west,font=\sffamily\scriptsize,text=accent] at (0,0) {directory = 1};
  \node[anchor=west,font=\sffamily\scriptsize,text=green]  at (3.4,0) {table = 6};
  \node[anchor=west,font=\sffamily\scriptsize,text=amber]  at (6.8,0) {offset = 0xA2C = 2604};
\end{tikzpicture}
\end{center}

Check it with arithmetic instead of bit-counting:
\begin{align*}
\text{directory} &= \hexa{00401A2C} \gg 22 = 4{,}201{,}004 \div 4{,}194{,}304 = 1\\
\text{table} &= (\hexa{00401A2C} \gg 12) \,\&\, \hexa{3FF} = 1030 \,\&\, 1023 = 6\\
\text{offset} &= \hexa{00401A2C} \,\&\, \hexa{FFF} = \hexa{A2C} = 2604
\end{align*}

\textbf{Step 3 --- the hardware walks the tables.}

\begin{enumerate}
  \item Read \reg{CR3}. Say it holds \hexa{00018000}. That is the physical address of this
        process's page directory.
  \item Go to entry \textbf{1} of that directory: physical address
        $\hexa{18000} + 1 \times 4 = \hexa{18004}$. Suppose it contains \hexa{0001B007}.
  \item Mask off the flags: $\hexa{0001B007} \,\&\, \hexa{FFFFF000} = \hexa{0001B000}$. That is the
        physical address of the page table.
  \item Go to entry \textbf{6} of that table: $\hexa{1B000} + 6 \times 4 = \hexa{1B018}$. Suppose it
        contains \hexa{00C4D007}.
  \item Mask again: \hexa{00C4D000}. That is the physical address of the \emph{frame}.
  \item Add the offset: $\hexa{00C4D000} + \hexa{A2C} = \hi{\hexa{00C4DA2C}}$.
\end{enumerate}

So virtual \hexa{00401A2C} is physical \hexa{00C4DA2C}. The program has no idea. It never will.

\begin{center}
\begin{tikzpicture}[font=\sffamily\scriptsize,
  n/.style={draw=grayborder,rounded corners=2pt,minimum width=2.7cm,minimum height=0.85cm,align=center}]
  \node[n,fill=graylight] (v) at (0,3.2) {virtual\\\texttt{0x00401A2C}};
  \node[n,fill=redlight,draw=red] (cr3) at (0,1.6) {CR3\\\texttt{0x00018000}};
  \node[n,fill=accentlight,draw=accent] (dir) at (4.2,1.6) {directory[1]\\\texttt{0x0001B000}};
  \node[n,fill=greenlight,draw=green] (tab) at (8.4,1.6) {table[6]\\\texttt{0x00C4D000}};
  \node[n,fill=amberlight,draw=amber] (fin) at (8.4,0) {+ offset 0xA2C\\\texttt{0x00C4DA2C}};

  \draw[-{Stealth[length=2mm]},gray] (v) -- node[left,font=\tiny]{split} (cr3);
  \draw[-{Stealth[length=2mm]},accent] (cr3) -- node[above,font=\tiny]{index 1} (dir);
  \draw[-{Stealth[length=2mm]},green] (dir) -- node[above,font=\tiny]{index 6} (tab);
  \draw[-{Stealth[length=2mm]},amber] (tab) -- (fin);
\end{tikzpicture}
\end{center}

\begin{concept}{This happens on every single access}
Not once per program. Not once per function. \hi{Every time the processor reads or writes
memory} --- every instruction fetch, every variable access, every stack push --- it performs this
two-step lookup.

If it really did two extra memory reads per access, computers would run three times slower. They do
not, because of a cache called the \term{TLB} (Translation Lookaside Buffer) that remembers recent
translations. The TLB is why the next section matters.
\end{concept}

\subsection{A second example: a translation that fails}

The first example worked. Now let us take an address in the \emph{same} address space that does not
resolve, because understanding failure is what makes the mechanism click.

We will translate \hexa{00500000} --- five megabytes in --- inside the shell's address space.

\textbf{Step 1 --- split it.}

\begin{align*}
\text{directory} &= \hexa{00500000} \gg 22 = 1\\
\text{table} &= (\hexa{00500000} \gg 12) \,\&\, \hexa{3FF} = 1280 \,\&\, 1023 = 256\\
\text{offset} &= \hexa{00500000} \,\&\, \hexa{FFF} = 0
\end{align*}

\textbf{Step 2 --- the directory lookup succeeds.}

Directory entry 1 is present: it covers the whole 4 MB region \hexa{00400000}--\hexa{007FFFFF}, and
\hexa{500000} falls inside it. So far, so good. The hardware reads the entry, masks off the flags,
and walks to the page table.

\textbf{Step 3 --- the table lookup fails.}

Entry 256 of that table is \hi{zero}. Bit 0 --- the present bit --- is clear. The address has
a table to look in, and the table has nothing to say about it.

The hardware stops here. It does not guess, it does not return zero, it does not read whatever
happens to be at physical \hexa{500000}. It raises \hi{exception 14}, and before jumping to
the handler it leaves two pieces of evidence:

\begin{center}
\begin{tabularx}{0.93\textwidth}{@{}l l X@{}}
\toprule
\textbf{Where} & \textbf{Value} & \textbf{Meaning} \\
\midrule
\reg{CR2} & \hexa{00500000} & The address that was being translated \\
error code & \hexa{4} or \hexa{6} & bit 0 clear: \emph{not present}; bit 2 set: from ring 3; bit 1 set if it was a write \\
\bottomrule
\end{tabularx}
\end{center}

So a user program reading that address produces error code \hexa{4}, and writing to it produces
\hexa{6}. Run it through the handler from Chapter 28 and you get:

\begin{lstlisting}[style=txt]
PAGE FAULT
Address: 0x00500000
Error code: 0x00000006
Reason: Page not present
Access: Write
Mode: User
\end{lstlisting}

\begin{concept}{Two ways to fail, and they mean different things}
Notice there are two distinct failures hiding here, and the error code distinguishes them:

\begin{itemize}
  \item \hi{Bit 0 clear --- not present.} There is no translation. The name means nothing in
        this address space. This is a wild pointer, a stack that grew too far, or memory that was
        never allocated.
  \item \hi{Bit 0 set --- protection violation.} There \emph{is} a translation, but the access
        was not allowed: writing to a read-only page, or a ring 3 program touching a page whose
        \cod{PAGE\_USER} bit is clear. The page exists; you just may not do that to it.
\end{itemize}

The second kind is how the kernel stays protected. When a user program tries to read kernel memory,
the translation succeeds --- the kernel \emph{is} mapped into its address space --- and it is the
permission check that stops it, producing error code \hexa{5} (present + user). A program cannot
tell the difference between ``kernel memory'' and ``memory that does not exist'' except by the error
code it never gets to see.
\end{concept}

\begin{warn}{Why an unmapped address is not always an error}
In a full operating system, most page faults are not bugs. The page might be on disk (swapped out),
or it might be a stack page the kernel deliberately did not allocate until first touched, or a
copy-on-write page being written for the first time. In all of those cases the handler fixes the
situation, and then \emph{re-executes the faulting instruction}, which now succeeds. The program
never finds out.

lytlnyblOS does none of this: every fault is fatal to the process that caused it. But the mechanism
is identical, and Part XIII lists on-demand stack growth as one of the most rewarding extensions to
attempt --- it is about fifteen lines, and it turns this exact failure into a feature.
\end{warn}

\section{Why a page is 4096 bytes}

The number looks arbitrary. It is not, and seeing why fixes the whole design in your mind.

Page size is a three-way trade-off, and 4 KB is close to the sweet spot for a 32-bit machine.

\subsection{The cost of a small page}

Halve the page to 2 KB and you need twice as many entries to describe the same memory. Worse, the
arithmetic stops being neat: 2048 bytes needs 11 offset bits, leaving 21 to split between two levels
--- which do not divide evenly, and whichever way you cut it the tables no longer fit in one page.

Take it down to 1 KB and a two-level split gives 11 and 11 bits: tables of 2048 entries, 8192 bytes
each. A page table now spans \emph{eight} pages, which the allocator has to find contiguously. The
elegance is gone.

\subsection{The cost of a large page}

Go up to 16 KB and the tables shrink to 512 entries --- 2 KB, wasting most of a 16 KB page each. And
now \term{internal fragmentation} bites: every allocation wastes on average half a page.

\begin{center}
\begin{tabularx}{\textwidth}{@{}l l l l X@{}}
\toprule
\textbf{Page} & \textbf{Offset bits} & \textbf{Split} & \textbf{Table size} & \textbf{Average waste per allocation} \\
\midrule
1 KB   & 10 & 11 + 11 & 8 KB (8 pages!) & 512 B \\
\textbf{4 KB} & \textbf{12} & \textbf{10 + 10} & \textbf{4 KB (exactly 1 page)} & \textbf{2 KB} \\
16 KB  & 14 & 9 + 9  & 2 KB (1/8 of a page) & 8 KB \\
\bottomrule
\end{tabularx}
\end{center}

\begin{concept}{4 KB is the size where the design closes on itself}
Only at 4096 bytes does everything line up:
\begin{itemize}
  \item 12 offset bits leave exactly 20, which split evenly into 10 + 10;
  \item $2^{10} = 1024$ entries, at 4 bytes each, is exactly 4096 bytes;
  \item so \hi{a page table is exactly one page}.
\end{itemize}
That last point is what makes the implementation so short. \cod{create\_page\_table} can just call
\cod{alloc\_frame()} and use the result directly --- no special allocator, no contiguity problem, no
leftover space. Look at how little code it takes:
\begin{lstlisting}[style=c]
page_table_t* create_page_table(page_directory_t* directory,
                                uint32_t directory_index, uint32_t flags) {
    page_table_t* new_table = alloc_frame();
    memset(new_table, 0, 4096);
    (*directory)[directory_index] = (uintptr_t)new_table | flags;
    return new_table;
}
\end{lstlisting}
Four lines, because the page size was chosen so that four lines would do.
\end{concept}

\begin{warn}{And the reason modern systems also offer huge pages}
Each TLB entry covers one page. With 4 KB pages, a program touching 1 GB of memory needs 262,144
distinct translations, and the TLB holds a few hundred --- so it misses constantly.

This is why x86 also supports 2 MB and 1 GB pages. One TLB entry then covers 2 MB, and the same
program needs 512 translations instead of 262,144. The cost is fragmentation: allocate a 2 MB page
for a 4 KB structure and you have wasted 2,093,056 bytes.

Databases and virtual machines use huge pages for exactly this reason. lytlnyblOS does not --- the
bit that enables them lives in the directory entry, and turning it on is a small but genuinely
instructive change.
\end{warn}

\section{The TLB and why stale entries kill}

The TLB caches virtual $\rightarrow$ physical translations. It makes paging affordable. It also
means that \hi{if you change a page table, the processor may not notice}, because it is still
using a cached answer.

Two instructions fix that:

\filetag{lytlnybl/memory/vmm.asm}
\begin{lstlisting}[style=asm]
flush_tlb:              ; throw away EVERYTHING
    mov eax, cr3
    mov cr3, eax        ; reloading CR3 with itself clears the whole TLB
    ret

flush_tlb_page:         ; throw away ONE entry
    mov eax, [esp + 4]
    invlpg [eax]
    ret
\end{lstlisting}

\cod{invlpg} invalidates a single page's cached translation, which is far cheaper than dumping the
whole TLB. The mapping code uses it:

\filetag{lytlnybl/memory/vmm.c}
\begin{lstlisting}[style=c]
(*selected_table)[table_index] = physical_address | flags;
if (directory == current_directory) {
    flush_tlb_page(virtual_address);
}
\end{lstlisting}

\begin{warn}{Note the condition}
The flush only happens if we are modifying the \emph{currently active} directory. Modifying another
process's directory needs no flush, because the TLB holds nothing from it --- and when that process
is scheduled, \reg{CR3} is reloaded, which clears the TLB anyway.

Forgetting a flush produces one of the most maddening bugs in systems programming: you map a page,
the table clearly contains the right entry, you can print it and see it --- and the access still
faults, because the processor is answering from a cache you cannot see.
\end{warn}

\section{Anatomy of a table entry}

Each entry is 32 bits. The top 20 hold a physical address (which is enough, because everything is
page-aligned so the low 12 bits are always zero). Those freed-up 12 bits carry the flags:

\begin{center}
\begin{tikzpicture}[font=\sffamily\scriptsize,x=1cm]
  \draw[draw=grayborder,fill=accentlight] (0,0) rectangle (7.2,0.95);
  \node at (3.6,0.47) {physical address, bits 31--12};
  \foreach \i/\n/\c in {0/G/graylight,1/D/amberlight,2/A/amberlight,3/{}/graylight,4/{}/graylight,
                        5/U/greenlight,6/W/greenlight,7/P/redlight} {
    \draw[draw=grayborder,fill=\c] (7.2+\i*0.7,0) rectangle (7.9+\i*0.7,0.95);
    \node[font=\ttfamily\scriptsize] at (7.55+\i*0.7,0.47) {\n};
  }
  \node[font=\ttfamily\tiny,text=gray] at (0.5,-0.3) {bit 31};
  \node[font=\ttfamily\tiny,text=gray] at (12.4,-0.3) {bit 0};
\end{tikzpicture}
\end{center}

\begin{center}
\begin{tabularx}{\textwidth}{@{}c l l X@{}}
\toprule
\textbf{Bit} & \textbf{Name} & \textbf{Constant} & \textbf{Meaning} \\
\midrule
0 & P & \cod{PAGE\_PRESENT} & \textbf{The page exists.} If 0, any access raises a page fault \\
1 & W & \cod{PAGE\_WRITABLE} & Writable. If 0, the page is read-only \\
2 & U & \cod{PAGE\_USER} & \textbf{Ring 3 may access it.} If 0, kernel only \\
5 & A & --- & Accessed; set by the processor on any access \\
6 & D & --- & Dirty; set by the processor on a write \\
\bottomrule
\end{tabularx}
\end{center}

\filetag{lytlnybl/include/memory/vmm.h}
\begin{lstlisting}[style=c]
#define PAGE_PRESENT  0x1
#define PAGE_WRITABLE 0x2
#define PAGE_USER     0x4
\end{lstlisting}

\begin{concept}{Bit 2 is where user isolation actually lives}
Everything in Part X about user space comes down to this bit. A kernel page has \cod{PAGE\_USER}
clear; a user page has it set. When a ring 3 program touches an address whose page lacks that bit,
the processor raises a page fault \emph{before} the access happens. No data leaks. No check in
software. One bit, enforced in silicon, on every access.

And the corollary: if you accidentally set \cod{PAGE\_USER} on a kernel page, you have handed user
programs the keys, and nothing will warn you.
\end{concept}

\begin{warn}{Permissions are checked at both levels}
The directory entry \emph{also} has W and U bits, and the effective permission is the
\emph{stricter} of the two. A page marked writable in the table but read-only in the directory is
read-only. Forgetting this leads to faults that make no sense until you remember to look one level
up.
\end{warn}

\section{A real address space, captured live}

Everything so far has been described. This section shows the actual bytes, read out of a running
lytlnyblOS while the shell sat at its prompt.

\begin{tryit}[How this was captured]
QEMU's monitor can read guest memory. Started with
\cod{-monitor unix:/tmp/mon.sock,server,nowait}, the commands are \cod{memsave} (virtual addresses,
using the current CPU state) and \cod{pmemsave} (physical addresses).

The kernel's \cod{process\_head} global is at \hexa{13400} in this build --- find it with
\cod{i686-elf-nm build/kernel.elf}. Reading four bytes there gives the head of the process list;
walking \cod{next} at offset 96 and reading \cod{page\_directory} at offset 100 gives each process's
directory, which \cod{pmemsave} can then dump in full.
\end{tryit}

\subsection{The process list}

\begin{lstlisting}[style=txt]
process_head -> 0xc000000c
  pid=1  KERNEL RUNNING  pd=0x00001000 cs=0x08 eip=0x0000b022 esp=0x00012998
  pid=2  KERNEL RUNNING  pd=0x00001000 cs=0x08 eip=0x0000ac08 esp=0x00012534
  pid=3  USER   BLOCKED  pd=0x00019000 cs=0x1b eip=0x004010f2 esp=0xbfffef9c
\end{lstlisting}

Read that carefully, because several things from earlier chapters are visible at once:

\begin{itemize}
  \item \cod{process\_head} is \hexa{C000000C} --- in the kernel heap, twelve bytes past
        \hexa{C0000000}, which is exactly \cod{sizeof(heap\_header\_t)}. It is the very first
        \cod{kmalloc} the system ever performs.
  \item Both kernel processes share \cod{pd=0x00001000}: the one kernel directory. The shell has its
        own at \hexa{19000}.
  \item The shell's \cod{cs} is \hexa{1B}. That is \hexa{18}\cod{ | 3} --- the user code selector
        with RPL 3. \hi{This is the number that proves it is running unprivileged.}
  \item Its \cod{eip} is \hexa{4010F2}: inside the code mapped at \hexa{400000}. Its \cod{esp} is
        \hexa{BFFFEF9C}: just under \hexa{BFFFF000}, the stack top.
  \item Its state is \cod{BLOCKED} --- it is waiting in \cod{read()} for a keystroke, exactly as
        Chapter 24 described.
\end{itemize}

\subsection{The shell's page directory}

Dumping all 4096 bytes at physical \hexa{19000} and listing only the non-zero entries gives this ---
\hi{five entries out of 1024}:

\begin{lstlisting}[style=txt]
dir[  0] = 0x00002023  table@0x00002000 flags=0x023  covers 0x00000000-0x003fffff
dir[  1] = 0x0001f027  table@0x0001f000 flags=0x027  covers 0x00400000-0x007fffff
dir[  2] = 0x0001d007  table@0x0001d000 flags=0x007  covers 0x00800000-0x00bfffff
dir[767] = 0x0001b027  table@0x0001b000 flags=0x027  covers 0xbfc00000-0xbfffffff
dir[768] = 0x00004023  table@0x00004000 flags=0x023  covers 0xc0000000-0xc03fffff
\end{lstlisting}

This one listing is the entire argument for two-level tables. A flat table would need 1,048,576
entries and 4 MB. This process describes its whole world with a 4 KB directory and four 4 KB tables:
\hi{20 KB instead of 4 MB}, because 1019 of the 1024 entries are simply zero.

Decoding the flags with the table from the previous section:

\begin{center}
\begin{tabularx}{\textwidth}{@{}l l l X@{}}
\toprule
\textbf{Entry} & \textbf{Flags} & \textbf{Decoded} & \textbf{What it is} \\
\midrule
dir[0]   & \hexa{023} & P, W, A       & The identity-mapped kernel. \textbf{No U bit} \\
dir[1]   & \hexa{027} & P, W, U, A    & User code at \hexa{400000} \\
dir[2]   & \hexa{007} & P, W, U       & User heap region. \textbf{No A bit --- never touched yet} \\
dir[767] & \hexa{027} & P, W, U, A    & User stack \\
dir[768] & \hexa{023} & P, W, A       & The kernel heap. \textbf{No U bit} \\
\bottomrule
\end{tabularx}
\end{center}

\begin{concept}{The U bit, doing its job, visible in real data}
Three of these regions are reachable from ring 3 and two are not, and the only difference is bit 2.
The kernel is \emph{present} in this address space --- it has to be, or interrupts could not be
handled --- but \cod{dir[0]} and \cod{dir[768]} lack the U bit, so the shell cannot read a single
byte of it. Not because anything checks: because the hardware refuses.

That is the whole of user-space isolation, and here it is as two hexadecimal digits differing by 4.
\end{concept}

\begin{warn}{The accessed bit tells you what has actually run}
\cod{dir[2]} has no A bit, while every other entry does. The processor sets A on first access, so
this says the shell \emph{has never touched its heap}.

And that is exactly right: the shell blocks in \cod{read()} on its very first loop iteration, before
it ever reaches \cod{tokenize\_line}, which is where the first \cod{calloc} happens. The hardware is
reporting, accurately, that a line of C has not been reached yet. It is a small thing, but it is a
good demonstration that these bits are real and live.
\end{warn}

\subsection{Inside the tables}

Dumping the four tables and listing the non-zero entries:

\begin{lstlisting}[style=txt]
table for dir[1] (0x00400000) : 2 entries
   [   0] 0x0001e027 -> phys 0x0001e000  [P,W,U,A]   virt 0x00400000
   [   1] 0x00020027 -> phys 0x00020000  [P,W,U,A]   virt 0x00401000

table for dir[2] (0x00800000) : 1 entry
   [ 512] 0x0001c007 -> phys 0x0001c000  [P,W,U]     virt 0x00a00000

table for dir[767] (0xbfc00000) : 1 entry
   [1022] 0x0001a067 -> phys 0x0001a000  [P,W,U,A,D] virt 0xbfffe000

table for dir[0] (0x00000000) : 1024 entries, contiguous identity map
   [   0] 0x00000003 -> phys 0x00000000  [P,W]       virt 0x00000000
   [   1] 0x00001063 -> phys 0x00001000  [P,W,A,D]   virt 0x00001000
   ...
   [1023] 0x003ff003 -> phys 0x003ff000  [P,W]       virt 0x003ff000
\end{lstlisting}

Every number here can be predicted from the code, which is the best possible confirmation that you
have understood it:

\begin{description}
  \item[Two code pages, not one] \cod{rootfs/bin/shell} is 4,392 bytes, and
    $\lceil 4392 / 4096 \rceil = 2$. That is precisely the loop bound in \cod{load\_program}. The
    two frames, \hexa{1E000} and \hexa{20000}, are \emph{not adjacent} --- \hexa{1F000} went to the
    page table itself, allocated in between. Physical scattering, virtual contiguity: the entire
    point of paging, visible in three numbers.
  \item[The heap page is at table index 512] \cod{USER\_HEAP\_START} is \hexa{A00000}.
    $(\hexa{A00000} \gg 12) \,\&\, \hexa{3FF} = 512$. It is the only mapped page in that 4 MB
    region, which is why the region shows one entry out of 1024.
  \item[The stack page is at index 1022] \cod{USER\_STACK\_TOP} is \hexa{BFFFF000} and
    \cod{create\_uprocess} maps \cod{USER\_STACK\_TOP - 4096} = \hexa{BFFFE000}.
    $(\hexa{BFFFE000} \gg 12) \,\&\, \hexa{3FF} = 1022$. And the saved \cod{esp} was \hexa{BFFFEF9C},
    which is inside that page --- 3,996 bytes into it, with 100 bytes of stack used.
  \item[The stack page is the only dirty one] Flags \hexa{067} include D. The processor sets D on
    the first \emph{write}. Code pages are read and executed but not written; the heap has not been
    touched; only the stack has been written to. The hardware has independently confirmed what the
    program does.
\end{description}

\begin{danger}{And the one-page stack limit, in the same listing}
That table has exactly one entry. The next page down, \hexa{BFFFD000} at index 1021, is zero.

So the moment \reg{ESP} drops below \hexa{BFFFE000} --- 4,096 bytes of stack --- the translation
fails exactly like the \hexa{500000} example earlier, and the process dies. The saved \cod{esp}
shows 100 bytes in use, so there are 3,996 bytes of headroom. A local array of 4 KB, or roughly
forty nested calls, and the shell is gone.

This is the limitation described in Chapter 33, and here is the measurement that proves it.
\end{danger}

\section{The same address in two address spaces}

The promise at the start of this chapter was that two programs could both use \hexa{400000} without
colliding. The dumps let us check the claim directly.

In the \emph{shell's} directory, entry 1 is present and maps \hexa{400000} to physical
\hexa{1E000}. In the \emph{kernel's} directory --- dumped separately from physical \hexa{1000} ---
entry 1 is:

\begin{lstlisting}[style=txt]
dir[0] = 0x00002023   (the identity map)
dir[1] = 0x00000000   <- nothing
\end{lstlisting}

\begin{concept}{The same number, two different truths}
\hexa{400000} is a perfectly valid address in the shell's world: it holds the first instruction the
shell executes. In the kernel's world the identical number means nothing at all --- touching it
raises a page fault.

Nothing about the address changed. What changed is which dictionary was open, and \reg{CR3} is what
selects the dictionary. \hi{Changing \reg{CR3} changes what every address in the machine
means}, in one instruction.

That is why \cod{set\_cr3} appears in \cod{context\_switch}, and why the keyboard handler has to
change it before writing into a sleeping process's buffer. And it is why running a second copy of a
program needs no relocation, no cooperation and no trust: the loader allocates fresh frames, writes
them into \emph{that} process's table at the same virtual address, and the two programs are now
mutually invisible while both believing they live at \hexa{400000}.
\end{concept}

\begin{warn}{With one deliberate exception}
Look again at \cod{dir[0]} in both directories: \hi{both are \hexa{00002023}, pointing at the
same page table at physical \hexa{2000}.}

That is not a coincidence and not a copy --- it is literally the same table, shared. It is what
\cod{create\_uprocess} does when it copies all 1024 directory entries from the kernel directory, and
it is why the kernel appears at the same addresses in every process.

It is also why the bug in \cod{destroy\_process} described in Chapter 32 is serious: calling
\cod{unmap\_page} on \cod{kernel\_directory} edits that shared table, so the hole it punches appears
in every address space in the system at once.
\end{warn}

\begin{recipe}{mapping, aliasing, translating and unmapping}{ch27-paging}
The whole of this chapter as running code. It maps a frame at \hexa{40000000} --- chosen because it
is above the identity map, below the kernel heap and outside every user region --- writes through
the virtual address and reads the value back through the physical one, then maps a \emph{second}
virtual address onto the same frame and shows a write through one appearing through the other.

It then splits the address by hand, compares against \cod{get\_physical\_address}, unmaps, and
builds a complete new address space with the kernel copied in. Along the way it demonstrates that
\cod{get\_physical\_address} never checks the present bit, so it keeps reporting an address for a
page you have just unmapped.
\end{recipe}


\exercise{The virtual address \hexa{C0000000} is the start of the kernel heap. Split it into its
three fields by hand and state which directory entry it uses. (Check: the answer is a well-known
round number, and it explains why \hexa{C0000000} was chosen rather than, say, \hexa{C0000400}.)}

% ============================================================
\chapter{Paging in practice}

\begin{goals}
  \item Read the real mapping code line by line
  \item Understand identity mapping and why the kernel needs it
  \item See how a process gets its own address space
\end{goals}

\section{Turning it on}

\filetag{lytlnybl/memory/vmm.c}
\begin{lstlisting}[style=c]
void init_vmm() {
    kernel_directory   = alloc_frame();
    page_table_t* tbl_ptr = alloc_frame();

    memset(kernel_directory, 0, 4096);
    memset(tbl_ptr, 0, 4096);

    /* map the first 4 MB with identity: virtual X -> physical X */
    for (uint32_t i = 0; i < 1024; i++) {
        (*tbl_ptr)[i] = (i * 0x1000) | PAGE_PRESENT | PAGE_WRITABLE;
    }

    (*kernel_directory)[0] = ((uintptr_t)tbl_ptr) | PAGE_PRESENT | PAGE_WRITABLE;

    current_directory = kernel_directory;
    set_cr3((uintptr_t)kernel_directory);
    set_cr0();      /* enable paging: the moment of truth */
}
\end{lstlisting}

The loop fills entry $i$ with physical address $i \times 4096$. So entry 0 maps virtual
\hexa{0000}--\hexa{0FFF} to physical \hexa{0000}--\hexa{0FFF}, entry 1 maps \hexa{1000} to
\hexa{1000}, and so on for 1024 entries: exactly 4 MB.

\begin{concept}{Why identity mapping is necessary, not just convenient}
Think about what happens the instant \cod{set\_cr0()} executes. Before it, addresses were physical.
After it, they are virtual. But \reg{EIP} does not change --- the very next instruction is fetched
from the same numeric address it would have been anyway.

If that address did not map to itself, the processor would fetch the next instruction from
somewhere completely different, and the system would die instantly and inexplicably.

Identity mapping the region containing the kernel makes the transition a non-event. The kernel sits
at \hexa{9000}, the stack, the bitmap at \hexa{100000} and the video buffer at \hexa{B8000} are all
below 4 MB, so all of them keep working unchanged.
\end{concept}

\begin{danger}{Only 4 MB is identity-mapped}
\cod{alloc\_frame()} hands out frames from anywhere in physical RAM --- but the kernel can only
\emph{reach} the first 4 MB. If an allocation ever returns a frame above \hexa{400000}, the
\cod{memset} in \cod{create\_page\_table} or the \cod{memcpy} in the program loader will touch an
unmapped address and fault.

It does not happen today because \cod{alloc\_frame} scans from frame 0 and the system uses only a
few dozen frames. But it is a ceiling, not a safety property. Any change that increases memory
usage --- more processes, a bigger heap, larger programs --- walks towards it. The proper fixes are
either to identity-map more of RAM, or to give the kernel a temporary mapping window it can point
at any frame.
\end{danger}

\section{Mapping a page}

\filetag{lytlnybl/memory/vmm.c}
\begin{lstlisting}[style=c]
void map_page(page_directory_t* directory, uintptr_t virtual_address,
              uintptr_t physical_address, uint32_t flags) {
    uint16_t dir_index   = (virtual_address >> 22);
    uint16_t table_index = (virtual_address >> 12 & 0x3FF);

    uint32_t dir_entry = (*directory)[dir_index];
    page_table_t* selected_table;

    if (!(dir_entry & PAGE_PRESENT)) {
        /* no table for this 4 MB region yet: create one */
        selected_table = create_page_table(directory, dir_index, flags);
    } else {
        selected_table = (page_table_t*)(dir_entry & 0xFFFFF000);
    }

    (*selected_table)[table_index] = physical_address | flags;

    if (directory == current_directory) {
        flush_tlb_page(virtual_address);
    }
}
\end{lstlisting}

This is the function that does everything. Split the address, find or create the table, write the
entry, flush if needed.

\begin{warn}{A subtlety in \cod{create\_page\_table}}
Notice that the same \cod{flags} are used for both the page table entry and the directory entry. If
you map a user page with \cod{PAGE\_USER}, the newly created directory entry also gets
\cod{PAGE\_USER} --- which is what you want here. But it means the \emph{first} page mapped into a
4 MB region decides the directory-level permissions for that whole region. Map a kernel-only page
first and a user page second into the same region, and the user page will be unreachable from ring
3 because the directory entry lacks the U bit. This is a real trap.
\end{warn}

\section{A process's own address space}

\filetag{lytlnybl/tasks/procman.c}
\begin{lstlisting}[style=c]
void create_uprocess(process_t* new_process) {
    new_process->page_directory = (page_directory_t*)alloc_frame();
    memset((page_directory_t*)new_process->page_directory, 0, 4096);

    /* copy every kernel mapping */
    for (uint32_t i = 0; i < 1024; i++) {
        (*new_process->page_directory)[i] = (*kernel_directory)[i];
    }

    /* then add what is private to this process */
    uintptr_t ustack_frame = (uintptr_t)alloc_frame();
    map_page(new_process->page_directory, USER_STACK_TOP - 4096, ustack_frame,
             PAGE_PRESENT | PAGE_WRITABLE | PAGE_USER);

    map_page(new_process->page_directory, USER_HEAP_START,
             (uintptr_t)alloc_frame(),
             PAGE_PRESENT | PAGE_WRITABLE | PAGE_USER);
    /* ... */
}
\end{lstlisting}

\begin{concept}{Why the kernel is mapped into every process}
It looks wasteful and even dangerous to copy the kernel's mappings into every user address space.
It is neither, and it is what every real operating system does.

The reason is the interrupt. When a user process makes a system call or the timer fires, the
processor switches to ring 0 --- but \hi{it does not change \reg{CR3}}. The kernel therefore
begins executing with the \emph{user's} page directory still active. If the kernel's code were not
mapped there, the very first instruction of the handler would fault.

Safety comes from the \cod{PAGE\_USER} bit: those kernel mappings lack it, so they are present but
unreachable from ring 3. The kernel is visible to the hardware and invisible to the program.
\end{concept}

\begin{warn}{The tables are shared, not copied}
The copy loop copies \emph{directory entries}, which are pointers. Every process therefore points
at the \emph{same} kernel page tables. That is intentional and efficient --- a change to a kernel
mapping is instantly visible everywhere.

But it also means that unmapping something from \cod{kernel\_directory} affects every process
simultaneously. There is a place in \cod{destroy\_process} that does exactly that, and Part XIII
explains why it is a bug.

The source code even carries a warning about this:
\begin{lstlisting}[style=c]
/* if you ever try to create more mappings and get a general protection fault
 * it is likely that you have overwritten an existing kernel mapping, as we copy
 * the kernel mappings to the user process for when we use our interrupts, just
 * be careful of this */
\end{lstlisting}
\end{warn}

\section{The address space, laid out}

\begin{center}
\begin{tikzpicture}[font=\sffamily\scriptsize,x=1cm,y=1cm]
  \def\w{7.6}
  \newcommand{\banda}[5]{%
    \fill[#4,draw=grayborder] (0,#1) rectangle (\w,#2);
    \node[align=center,font=\sffamily\scriptsize] at (\w/2,{(#1+#2)/2}) {#3};
    \node[anchor=west,font=\ttfamily\tiny,text=gray] at (\w+0.15,#1) {#5};
  }
  \banda{0}{1.0}{\textbf{Identity-mapped kernel}\\{\tiny code, stack, bitmap, VGA}}{redlight}{0x00000000}
  \banda{1.0}{1.9}{\textbf{User code}\\{\tiny loaded program}}{greenlight}{0x00400000}
  \banda{1.9}{2.7}{\textbf{User heap}\\{\tiny grows upward via sbrk}}{greenlight}{0x00A00000}
  \banda{2.7}{3.3}{User VGA {\tiny (reserved)}}{graylight}{0x00B00000}
  \banda{3.3}{4.3}{\textit{unmapped}}{white}{}
  \banda{4.3}{5.1}{\textbf{User stack}\\{\tiny grows downward}}{greenlight}{0xBFFFF000}
  \banda{5.1}{6.1}{\textbf{Kernel heap}\\{\tiny kmalloc, grows upward}}{redlight}{0xC0000000}
  \node[anchor=west,font=\ttfamily\tiny,text=gray] at (\w+0.15,6.1) {0xFFFFFFFF};
\end{tikzpicture}
\end{center}

Note the deliberate design: the user stack grows down from \hexa{BFFFF000} and the user heap grows
up from \hexa{A00000}. They approach each other with a large gap between. If they ever met, one
would silently corrupt the other --- and nothing in this system checks for that.

% ============================================================
\chapter{Page faults}

\begin{goals}
  \item What a page fault is and why it is not always an error
  \item How to decode the error code
  \item How to read a fault and find the cause
\end{goals}

\section{The most useful exception}

A \term{page fault} (exception 14) fires whenever the processor cannot complete a translation. In a
mature operating system that is often \emph{normal} and even desirable: it is how demand paging,
swapping, copy-on-write and lazy stack growth are implemented. The kernel fixes the situation and
resumes the instruction as if nothing happened.

lytlnyblOS does none of that. Here a page fault always means a bug. But its handler is still the
single most informative diagnostic in the system.

\section{Two sources of information}

\begin{description}
  \item[\reg{CR2}] The processor puts \emph{the address that was accessed} here. This is the first
    thing to look at.
  \item[The error code] Pushed on the stack, it describes \emph{what kind} of access failed.
\end{description}

\begin{center}
\begin{tabularx}{\textwidth}{@{}c l X@{}}
\toprule
\textbf{Bit} & \textbf{If 0} & \textbf{If 1} \\
\midrule
0 & The page was not present & Protection violation (page existed, access not allowed) \\
1 & It was a read & It was a write \\
2 & From the kernel (ring 0) & \textbf{From a user program (ring 3)} \\
3 & --- & A reserved bit was set in a table entry \\
4 & --- & It was an instruction fetch \\
\bottomrule
\end{tabularx}
\end{center}

\filetag{lytlnybl/memory/vmm.c}
\begin{lstlisting}[style=c]
void page_fault_handler(registers_t* registers) {
    uintptr_t address = get_cr2();
    uint32_t error_code = registers->error_code;

    vga_text_writeline(&terminal, "PAGE FAULT");

    vga_text_write(&terminal, "Address: ");
    vga_text_write_hex(&terminal, address);
    vga_text_writeline(&terminal, "");

    if (error_code & PRESENT_FAULT)
        vga_text_writeline(&terminal, "Reason: Protection violation");
    else
        vga_text_writeline(&terminal, "Reason: Page not present");

    if (error_code & WRITE_FAULT) vga_text_writeline(&terminal, "Access: Write");
    else                          vga_text_writeline(&terminal, "Access: Read");

    if (error_code & USER_FAULT)  vga_text_writeline(&terminal, "Mode: User");
    else                          vga_text_writeline(&terminal, "Mode: Kernel");

    vga_text_write(&terminal, "Instruction address: ");
    vga_text_write_hex(&terminal, registers->eip);
    vga_text_writeline(&terminal, "");
}
\end{lstlisting}

\section{Reading a fault}

The combination of address, reason and mode usually identifies the bug immediately:

\begin{center}
\begin{tabularx}{\textwidth}{@{}l l X@{}}
\toprule
\textbf{Address} & \textbf{Reason} & \textbf{Almost certainly} \\
\midrule
\hexa{00000000} or near & not present & A null pointer dereference \\
Just below a stack region & not present & \textbf{Stack overflow} \\
A kernel address, mode User & protection & A user program reaching where it should not \\
A page you just mapped & not present & A missing TLB flush \\
An address in a mapped page & protection & Writing to a read-only page, or missing \cod{PAGE\_USER} \\
\bottomrule
\end{tabularx}
\end{center}

\begin{concept}{The stack overflow signature}
This one is worth memorising because it is the most common serious fault in a kernel. You see a
fault at an address \emph{just below} the bottom of a stack region, with \cod{Mode: Kernel} and
usually \cod{Access: Write}, and the faulting instruction is a \cod{push} or a function prologue.

Example from this project: a fault at \hexa{BFFFFFFC} with \reg{ESP} = \hexa{C0000000}. Since
\hexa{C0000000} is the start of the kernel heap, a stack had grown downwards right off the bottom of
the heap's first page. That single observation identified the entire failure chain described in
Part XIII.
\end{concept}

% ============================================================
\begin{recipe}{reading a page fault report}{ch28-page-fault}
Three lines that write to \hexa{50000000}, which nothing maps. The comment above the write predicts
the report --- address, reason, access, mode, and an error code of \hexa{2} --- before you run it,
so you can check your understanding against the processor's own answer.
\end{recipe}

% ============================================================
\chapter{The kernel heap}

\begin{goals}
  \item Why pages are not enough and a byte allocator is needed
  \item How a linked-list heap works
  \item Understand splitting and merging of blocks
  \item Know the limitations of this implementation
\end{goals}

\section{Why pages are not enough}

The physical manager hands out 4 KB at a time. But a \cod{process\_t} structure is about 100 bytes.
Giving it a whole page wastes 97\,\% of it, and with a few dozen structures you have burned megabytes
on nothing.

The \term{heap} sits on top: it asks the page manager for pages, and hands out \emph{bytes} from
within them.

\section{The design: a linked list of blocks}

Every block of the heap is preceded by a small header:

\filetag{lytlnybl/include/memory/heap.h}
\begin{lstlisting}[style=c]
typedef struct heap_header {
    size_t size;                 /* usable bytes in this block */
    bool free;                   /* is it available? */
    struct heap_header* next;    /* the next block */
} heap_header_t;
\end{lstlisting}

\begin{center}
\begin{tikzpicture}[font=\sffamily\scriptsize,x=1cm]
  \draw[draw=grayborder,fill=amberlight] (0,0)   rectangle (1.1,1.0);
  \draw[draw=grayborder,fill=greenlight] (1.1,0) rectangle (3.6,1.0);
  \draw[draw=grayborder,fill=amberlight] (3.6,0) rectangle (4.7,1.0);
  \draw[draw=grayborder,fill=redlight]  (4.7,0) rectangle (7.2,1.0);
  \draw[draw=grayborder,fill=amberlight] (7.2,0) rectangle (8.3,1.0);
  \draw[draw=grayborder,fill=greenlight] (8.3,0) rectangle (10.8,1.0);
  \node[font=\sffamily\tiny] at (0.55,0.5) {hdr};
  \node[font=\sffamily\tiny,align=center] at (2.35,0.5) {free\\120 B};
  \node[font=\sffamily\tiny] at (4.15,0.5) {hdr};
  \node[font=\sffamily\tiny,align=center] at (5.95,0.5) {used\\64 B};
  \node[font=\sffamily\tiny] at (7.75,0.5) {hdr};
  \node[font=\sffamily\tiny,align=center] at (9.55,0.5) {free\\3800 B};
  \draw[-{Stealth[length=2mm]},accent] (0.55,1.1) to[bend left=35] (4.15,1.1);
  \draw[-{Stealth[length=2mm]},accent] (4.15,1.1) to[bend left=35] (7.75,1.1);
  \node[font=\sffamily\tiny,text=accent] at (5.9,2.0) {next pointers};
\end{tikzpicture}
\end{center}

The header lives \emph{immediately before} the memory it describes. That is why \cod{kmalloc}
returns a pointer past the header, and \cod{kfree} works backwards to find it:

\filetag{lytlnybl/memory/heap.c}
\begin{lstlisting}[style=c]
return (void*)((uintptr_t)block + sizeof(heap_header_t));   /* in kmalloc */

heap_header_t* get_header(void* ptr) {                      /* in kfree */
    return (heap_header_t*)((uintptr_t)ptr - sizeof(heap_header_t));
}
\end{lstlisting}

\begin{analogy}
Every block is a parcel with a label stuck to the front. The label says how big the parcel is,
whether it is empty, and where the next parcel is. When you are handed a parcel you are given the
address of the \emph{contents}; to read the label you step backwards by the label's width.

And if you scribble past the end of your parcel, you write on the next parcel's label --- which is
why heap overflows corrupt the allocator itself rather than just neighbouring data.
\end{analogy}

\section{Allocating}

\filetag{lytlnybl/memory/heap.c}
\begin{lstlisting}[style=c]
void* kmalloc(size_t requested_size) {
    if (!requested_size) return NULL;

    heap_header_t* block = find_free_block(requested_size);
    if (block) {
        split_block(block, requested_size);
        block->free = false;
        return (void*)((uintptr_t)block + sizeof(heap_header_t));
    } else {
        expand_heap(requested_size);
        return kmalloc(requested_size);     /* try again */
    }
}
\end{lstlisting}

\subsection{Splitting}

If a 4000-byte block is used to satisfy a 100-byte request, the remaining 3900 bytes should not be
lost. \cod{split\_block} cuts it in two:

\begin{lstlisting}[style=c]
void split_block(heap_header_t* block, size_t requested_size) {
    size_t space_remaining = (block->size - requested_size);
    if (space_remaining < sizeof(heap_header_t) + sizeof(uint8_t)) {
        return;   /* not worth splitting: the remainder could not even hold a header */
    }
    /* ... create a new header in the middle and link it in ... */
}
\end{lstlisting}

That guard matters: a remainder smaller than a header plus one byte would produce a block that can
never be used, and whose header would overlap the next block.

\subsection{Freeing and merging}

\begin{lstlisting}[style=c]
void merge_blocks(heap_header_t* block) {
    if (!block->next || !block->next->free) return;
    heap_header_t* block_to_merge = block->next;

    block->size += block_to_merge->size + sizeof(heap_header_t);
    block->next  = block_to_merge->next;
    return merge_blocks(block);    /* keep going while the next is free */
}

void kfree(void* ptr) {
    if (!ptr) return;
    heap_header_t* block_to_free = get_header(ptr);
    block_to_free->free = true;
    merge_blocks(block_to_free);
}
\end{lstlisting}

Merging matters because without it the heap fragments: you end up with plenty of free bytes, none
of them adjacent, and a 200-byte request fails surrounded by free memory.

\begin{warn}{It only merges forwards}
\cod{merge\_blocks} joins a block with the ones \emph{after} it. If you free a block whose
\emph{predecessor} is already free, they stay separate --- the list is singly linked, so there is no
way to look back.

The classic fix is a doubly linked list, or storing a copy of the size at the end of each block (a
``footer'') so the previous block's header can be found by arithmetic. As it stands, the heap
fragments more than it needs to.
\end{warn}

\section{Growing}

\begin{lstlisting}[style=c]
void expand_heap(size_t required_size) {
    uint32_t required_pages = (required_size + 4096 - 1) / 4096;
    for (size_t i = 0; i < required_pages; i++) {
        void* new_page_physical_start = alloc_frame();
        map_page(kernel_directory, heap_end,
                 (uintptr_t)new_page_physical_start,
                 PAGE_PRESENT | PAGE_WRITABLE);
        /* ... append or extend the last block ... */
        heap_end += 0x1000;
    }
}
\end{lstlisting}

There is the round-up idiom again, this time used correctly. And there is the dependency chain from
Chapter 6 made concrete: the heap calls \cod{map\_page} (virtual memory), which calls
\cod{alloc\_frame} (physical memory). Three layers, each built on the one below.

\begin{warn}{No alignment}
\cod{kmalloc(1)} followed by \cod{kmalloc(4)} returns a 4-byte block at an odd address. On x86
unaligned access works, just slower; on many other architectures it faults outright. A production
allocator rounds every request up to 8 or 16 bytes. Worth knowing if you ever port this.
\end{warn}

\begin{recipe}{kmalloc, kfree, and watching the block list}{ch29-heap}
Allocates two blocks, then prints the whole block list after every operation, so you can watch
splitting happen and then watch the forward-only merge leave a hole behind.

The first block it shows is not one of ours: it is the 112-byte \cod{process\_t} that
\cod{init\_procman} allocated during boot --- the very first \cod{kmalloc} the system performs.

It also demonstrates something easy to trip over when writing your own examples: the file keeps the
full initialisation sequence even though it only uses the heap, because \cod{idt\_init} enables
interrupts, the timer then fires, and the scheduler walks a process list that
\cod{init\_procman} has to have built first.
\end{recipe}

\begin{tryit}
The kernel exercises the heap during boot:
\begin{lstlisting}[style=c]
uint32_t* numbers = (uint32_t*)kmalloc(5 * sizeof(uint32_t));
for (int i = 0; i < 5; i++) numbers[i] = (i + 1) * 10;
/* ... print them ... */
kfree(numbers);
\end{lstlisting}
If you see \cod{Values: 10 20 30 40 50} on screen, the whole three-layer stack works: physical
manager, paging and heap.
\end{tryit}

\begin{summary}
Memory is solved: the kernel can find free RAM, give each program its own private view of it, and
hand out individual bytes. Everything needed for real processes is now in place. That is the next
part.
\end{summary}
